% !TeX root = ../main.tex

\chapter{实验结果与分析}

\section{总体结果与主要发现}

\begin{table}[htbp]
  \centering
  \caption{五个 CNN 模型上的主要结果}
  \label{tab:main-results}
  \resizebox{\textwidth}{!}{
  \begin{tabular}{lcccccccc}
    \toprule
    模型 & ID Acc (\%) & FedViM $k$ & ACT $k$ & 压缩率 & FedViM Near (\%) & ACT Near (\%) & FedViM Far (\%) & ACT Far (\%) \\
    \midrule
    DenseNet169 & 96.50 & 1000 & 99  & 90.1\% & 80.04 & 95.57 & 84.78 & 96.17 \\
    EfficientNetV2-S & 97.01 & 512 & 69  & 86.5\% & 96.40 & 95.77 & 97.52 & 96.51 \\
    MobileNetV3-Large & 96.16 & 512 & 89  & 82.6\% & 96.31 & 95.26 & 97.34 & 96.42 \\
    ResNet101 & 96.22 & 1000 & 143 & 85.7\% & 95.86 & 96.28 & 96.68 & 96.73 \\
    ResNet50 & 96.53 & 1000 & 141 & 85.9\% & 95.68 & 95.32 & 97.19 & 95.89 \\
    \midrule
    平均 & 96.48 & 804.8 & 108.2 & 86.2\% & 92.86 & 95.64 & 94.70 & 96.34 \\
    \bottomrule
  \end{tabular}}
\end{table}

\begin{figure}[htbp]
  \centering
  \includegraphics[width=\textwidth]{figure_2_subspace_compression.png}
  \caption{五个 CNN 模型上 fixed-k FedViM 与 ACT-FedViM 的主子空间维度比较}
  \label{fig:subspace-compression}
\end{figure}

表~\ref{tab:main-results} 给出了五个 CNN 模型上 FedViM 与 ACT-FedViM 的主要结果，图~\ref{fig:subspace-compression} 展示了相应的主子空间压缩情况。可以看到，ACT-FedViM 在五款模型上均显著降低了主子空间维度。相较于 fixed-k FedViM，ACT 选择的 $k$ 均落在 69 至 143 的范围内，平均维度由 804.8 降至 108.2，平均压缩率达到 86.2\%。这表明，ACT 在当前实验范围内稳定地缩小了 ViM 后处理所需的主子空间规模。

从平均结果看，ACT-FedViM 的 Near-OOD AUROC 和 Far-OOD AUROC 分别由 92.86\% 提升至 95.64\%、由 94.70\% 提升至 96.34\%。但这一平均增益并不在所有模型上均匀出现：从逐模型结果看，ACT 在 DenseNet169 和 ResNet101 上带来性能提升，而在其余三款模型上则出现了轻微衰减。

除检测性能外，主子空间维度的降低还具有明确的部署含义。ViM 后处理需要保存主子空间矩阵 $P \in \mathbb{R}^{D \times k}$，并在打分阶段完成投影与残差计算，因此更小的 $k$ 直接对应更低的存储开销与计算开销。因此，ACT-FedViM 的更稳定收益主要体现在主子空间压缩与后处理轻量化，而其 AUROC 增益则具有一定模型依赖性。

\section{与 MSP、Energy 基线的整体比较}

\begin{table}[htbp]
  \centering
  \caption{不同方法的平均 OOD 检测表现}
  \label{tab:method-average}
  \begin{tabular}{lcc}
    \toprule
    方法 & 平均 Near-OOD AUROC (\%) & 平均 Far-OOD AUROC (\%) \\
    \midrule
    MSP & 90.42 & 87.27 \\
    Energy & 83.34 & 80.89 \\
    FedViM & 92.86 & 94.70 \\
    ACT-FedViM & 95.64 & 96.34 \\
    \bottomrule
  \end{tabular}
\end{table}

\begin{figure}[htbp]
  \centering
  \includegraphics[width=\textwidth]{figure_1_method_comparison.png}
  \caption{五个 CNN 模型上 FedViM、ACT-FedViM、MSP 与 Energy 的 Near/Far-OOD AUROC 对比}
  \label{fig:method-comparison}
\end{figure}

为了说明联邦 ViM 系方法的整体有效性，表~\ref{tab:method-average} 给出了四类方法在五个模型上的平均表现，图~\ref{fig:method-comparison} 则给出了对应的可视化结果。可以看出，基于特征空间结构的 ViM 系方法整体上优于纯输出空间基线。以平均结果计，FedViM 相比 MSP 的 Near-OOD 与 Far-OOD AUROC 分别提高 2.44 和 7.43 个百分点，相比 Energy 分别提高 9.51 和 13.82 个百分点；ACT-FedViM 相比 MSP 的 Near-OOD 与 Far-OOD AUROC 分别提高 5.22 和 9.08 个百分点，相比 Energy 分别提高 12.30 和 15.46 个百分点。

值得注意的是，在当前细粒度浮游生物任务中，Energy 的平均表现甚至低于 MSP。这一现象表明，单纯依赖输出空间置信度或能量信息的方法在该任务上存在明显局限；相比之下，利用特征空间结构信息的 ViM 系方法更适合当前 OOD 检测场景。由此可见，联邦化 ViM 本身已经构成一条有效的后处理路线。

\section{代表模型分析}

\begin{figure}[htbp]
  \centering
  \makebox[\textwidth][c]{\includegraphics[width=1.06\textwidth]{figure_3_selected_models.png}}
  \caption{MobileNetV3-Large、ResNet101 与 DenseNet169 三个代表模型上的 Near/Far-OOD AUROC 对比}
  \label{fig:selected-models}
\end{figure}

为覆盖不同模型上 ACT 收益的代表性结果形态，本文选取 MobileNetV3-Large、ResNet101 和 DenseNet169 三个模型进行进一步分析。三者分别对应以下三种典型情形：MobileNetV3-Large 在大幅压缩主子空间的同时性能仅出现有限波动；ResNet101 在显著压缩维度的同时保持与 fixed-k 方法接近甚至略优的检测能力；DenseNet169 则表现出 fixed-k 与自适应选维之间较大的结果差异。图~\ref{fig:selected-models} 展示了三类代表模型的差异模式。

MobileNetV3-Large 的 fixed-k 为 512，而 ACT 选择的主子空间维度为 89，压缩率达到 82.6\%。在此基础上，Near-OOD AUROC 由 96.31\% 变化为 95.26\%，Far-OOD AUROC 由 97.34\% 变化为 96.42\%。这一结果表明，在较轻量的 CNN 模型上，ACT 的主要收益并不一定表现为进一步提升 AUROC，而更多表现为在维持较高绝对性能的同时显著减小主子空间规模。

ResNet101 的 fixed-k 为 1000，ACT 选择 143 维主子空间，压缩率为 85.7\%。在更小主子空间下，Near-OOD AUROC 由 95.86\% 提升至 96.28\%，Far-OOD AUROC 由 96.68\% 变化至 96.73\%。这一结果说明，在高维 CNN 模型上，ACT-FedViM 可以在显著缩小子空间规模的同时保持与 fixed-k 方法基本一致的检测能力，并取得轻微增益。

DenseNet169 的 fixed-k 为 1000，而 ACT 仅保留 99 维主子空间，压缩率达到 90.1\%。在此设定下，Near-OOD AUROC 由 80.04\% 提升至 95.57\%，Far-OOD AUROC 由 84.78\% 提升至 96.17\%。与其他模型相比，该模型上 fixed-k 与自适应选维之间的差异最为显著。

综合三类代表模型可以看出，ACT 并不能一致提升 OOD 检测的 AUROC，其价值在于提供了统计驱动的、更紧凑的选维方式，并在某些特定模型架构上缓解了 fixed-k 可能带来的失配风险。
